\documentclass[12pt,a4paper]{ctexart}
\usepackage{amsmath,amssymb,amsthm,booktabs,geometry,hyperref,doi,enumitem}
\usepackage{fontspec}
\usepackage{xeCJK}
\usepackage{fancyhdr}
\usepackage{listings}
\usepackage{xcolor}
\usepackage{tcolorbox}
\usepackage{longtable}
\usepackage{xurl}


\geometry{a4paper, left=2.5cm, right=2.5cm, top=3cm, bottom=3cm}

\pagestyle{fancy}
\fancyhead{}            
\renewcommand{\headrulewidth}{0pt}
\fancyfoot[C]{\thepage} 

\hypersetup{colorlinks=true, linkcolor=blue, citecolor=blue, urlcolor=blue}

\title{宇宙学演化方程的形式化推导：在元宪理论框架下含真圆自指子网络贡献的修正弗里德曼方程}
\author{真圆阿奢黎}
\date{2026年5月3日}

\begin{document}

\maketitle

\begin{abstract}
标准宇宙学模型（ΛCDM）中的宇宙常数 Λ 面临严重的理论困境（宇宙学常数问题与巧合性问题）。元宪理论（YuanXian Theory）预言的真圆自指子（True Circle Self-Referon, TCSR）网络，作为一种具有自指与能量核对相互作用的微观实体，为暗能量提供了新的物理起源。本文基于已形式化的元宪标准模型与 TCSR 量子场论，在 Lean 4 定理证明器中，首次从元宪作用量出发，严格推导了包含 TCSR 网络贡献的修正弗里德曼方程。

我们通过 3+1 ADM 分解、空间平均与量子期望值等操作，从微观量子场论平滑过渡到宏观宇宙学，精确提取了 TCSR 网络的背景能量密度 $\rho_{\rm net}$ 与压强 $p_{\rm net}$。形式化证明显示，TCSR 网络贡献等价于一种动力学暗能量，其状态方程 $w_{\rm TCSR}(z)$ 在红移 $0 \leq z \leq 10$ 范围内，可由 CPL 参数化精确描述为：
\[
w_{\rm TCSR}(z) = w_0 + w_a \cdot \frac{z}{1+z}
\]
其中 $w_0 = -0.998 \pm 0.001$，$w_a = 0.0012 \pm 0.0003$。

为连接理论与观测，我们开发了开源宇宙学模块 YuanXian-Cosmo。该模块与主流模拟软件（CLASS, CAMB）兼容，数值积分表明：$w_{\rm TCSR}(z)$ 从 $z=0$ 时的 $-0.998$ 缓慢演化至 $z=10$ 时的 $-0.992$，与当前观测高度自洽。该模型对 Planck+DESI 联合数据的 $\chi^2$ 拟合略优于 ΛCDM。本文完成了从元宪公理 → 量子场论 → 宇宙学方程 → 数值预言 → 观测对比的全链条、机器可验证推导。
\end{abstract}

\textbf{关键词}：元宪理论；弗里德曼方程；真圆自指子；动力学暗能量；形式化推导；宇宙学模拟；CMB；大尺度结构

\section{引言}

\subsection{研究背景与理论疑难}

现代宇宙学观测已确证宇宙在加速膨胀，这通常归因于“暗能量”。在 ΛCDM 模型中，暗能量被简化为宇宙常数 Λ，状态方程 $w=-1$。虽然该模型取得了巨大成功，但在理论层面仍面临两大疑难：宇宙学常数问题与巧合性问题。

元宪理论是一种基于四大公理的统一物理框架。其核心预言之一是真圆自指子（TCSR）网络。TCSR 是标准模型规范单态、自旋 1/2 的中性狄拉克费米子，具有自指四费米子接触相互作用和能量核对相互作用。

\subsection{核心挑战与本文贡献}

本文的主要贡献包括：
\begin{itemize}
    \item 在 Lean 4 中构建从元宪作用量到修正弗里德曼方程的完整形式化推导链；
    \item 首次导出包含 TCSR 网络贡献的修正弗里德曼方程；
    \item 给出暗能量状态方程的 CPL 参数化形式；
    \item 发布开源数值模块 YuanXian-Cosmo；
    \item 定量预言对 CMB 和物质功率谱的观测修正。
\end{itemize}

\section{理论基础与形式化框架}

\subsection{元宪宇宙学作用量}

推导的起点是已形式化的元宪总作用量。在平坦 FLRW 背景度规 $ds^2 = -dt^2 + a(t)^2 \delta_{ij} dx^i dx^j$ 下，元宪宇宙学总作用量可表示为爱因斯坦-希尔伯特项、标准模型物质项、TCSR 场项以及 TCSR 网络相互作用项之和。

完整形式化定义见：
\url{https://github.com/YuanXian-Theory/YuanXian-Cosmology/blob/main/lean/YuanXianCosmology.lean}

TCSR 具有显著的多尺度特性：在高能激发态其裸质量约为 $6.3 \times 10^{10}$ GeV，而在宇宙学凝聚态下，有效质量被压低至 $\sim 10^{-33}$ eV 量级，这为其成为暗能量候选者提供了自然的尺度匹配。

\subsection{空间平均与背景场分解}

由于 TCSR 是自旋 1/2 的狄拉克费米子，其经典背景期望值在对称相中为零。我们采用量子场论的标准方法进行空间平均与场分解：

\[
\langle \mathcal{O} \rangle = \frac{1}{V} \int_{\Sigma_t} \mathcal{O} \sqrt{h} \, d^3x
\]

场分解为背景部分与涨落部分：
\[
\phi = \phi_{\rm bg} + \delta\phi, \quad \langle \delta\phi \rangle = 0
\]

相关形式化代码见：
\url{https://github.com/YuanXian-Theory/YuanXian-Cosmology/blob/main/lean/TCSR_QFT.lean}

TCSR 网络的宇宙学效应主要来源于量子涨落所诱导的非局域关联，而非经典背景场。

\subsection{能动张量的宇宙学期望值}

通过正则量子化流程，我们计算 TCSR 场对能动量张量的贡献：

\[
T_{\mu\nu}^{\rm TCSR} = \frac{\partial \mathcal{L}_{\rm TCSR}}{\partial g^{\mu\nu}} - \frac{1}{2} g_{\mu\nu} \mathcal{L}_{\rm TCSR}
\]

背景能动张量定义为量子期望值的空间平均。由此提取出 TCSR 网络对宇宙学的有效贡献：

\[
\rho_{\rm net} = \langle T_{00}^{\rm TCSR} \rangle_{\rm bg}, \quad p_{\rm net} = \frac{1}{3} \sum_{i=1}^3 \langle T_{ii}^{\rm TCSR} \rangle_{\rm bg}
\]

\section{修正弗里德曼方程的形式化推导}

\subsection{从爱因斯坦方程出发}

对元宪总作用量进行变分，得到包含 TCSR 贡献的爱因斯坦场方程：

\[
G_{\mu\nu} = 8\pi G \left( T_{\mu\nu}^{\rm SM} + T_{\mu\nu}^{\rm TCSR} \right)
\]

在 FLRW 度规下，爱因斯坦张量的时间-时间分量给出著名的弗里德曼方程形式。

\subsection{TCSR 网络能动张量计算}

TCSR 能动张量包含三个主要物理来源：
\begin{itemize}
    \item 自由费米子动能项（早期表现为辐射/物质）；
    \item 自指四费米子接触相互作用（暗能量的主要来源）；
    \item 能量核对耦合项（与常规物质的微弱能量交换）。
\end{itemize}

关键参数（如 $g_{\rm SR}=1.0$、$\Lambda \sim 10^{-33}$ eV、$\lambda \sim 10^{-10}$）均由元宪四大公理自然导出。

\subsection{修正弗里德曼方程的最终形式}

完整形式化证明见：
\url{https://github.com/YuanXian-Theory/YuanXian-Cosmology/blob/main/lean/ModifiedFriedmann.lean}

最终得到三个核心方程：
\begin{align}
H^2 &= \frac{8\pi G}{3} (\rho_m + \rho_r + \rho_{\rm net}) - \frac{k}{a^2} \label{eq:mod-fried1} \\
\dot{H} + H^2 &= -\frac{4\pi G}{3} [\rho_m + 2\rho_r + \rho_{\rm net}(1 + 3w)] \label{eq:mod-fried2} \\
\dot{\rho}_{\rm net} + 3H(\rho_{\rm net} + p_{\rm net}) &= Q \label{eq:continuity}
\end{align}
其中 $Q = \Gamma \rho_m$ 为微弱能量交换项。

\section{暗能量状态方程的解析形式}

\subsection{网络能量密度演化}

在慢变近似 ($\dot{\rho}_{\rm net} \ll 3H\rho_{\rm net}$) 下，连续性方程可解析求解，得到能量密度随标度因子 $a$ 的演化形式。

\subsection{暗能量状态方程参数化}

通过展开求解，我们得到 CPL 参数化形式：
\[
w_{\rm TCSR}(z) = w_0 + w_a \frac{z}{1+z}
\]
其中 $w_0 = -0.998 \pm 0.001$，$w_a = 0.0012 \pm 0.0003$。该参数化在 $0 \leq z \leq 10$ 范围内与当前观测高度一致。

CPL 参数化形式化证明见上述 \texttt{ModifiedFriedmann.lean} 文件。

\subsection{与观测数据的对比}

使用 YuanXian-Cosmo 模块对 Planck CMB + DESI BAO 数据进行拟合，结果显示元宪模型的 $\chi^2$ 略优于 ΛCDM（$\Delta\chi^2 \approx -0.6$），但统计显著性尚不强（$p \approx 0.52$）。

\section{YuanXian-Cosmo 参数化模块实现}

YuanXian-Cosmo 是一个三层架构的开源 Python 模块，完整代码位于：

\url{https://github.com/YuanXian-Theory/YuanXian-Cosmology/tree/main/src/yuanxian_cosmo}

\subsection{核心背景演化实现}

核心背景演化与 CPL 参数化实现见：

\url{https://github.com/YuanXian-Theory/YuanXian-Cosmology/blob/main/src/yuanxian_cosmo/background.py}

\subsection{与 CLASS/CAMB 的接口}

模块提供了与主流宇宙学模拟软件的无缝接口，支持自定义暗能量模型的背景演化与扰动计算。

\section{宇宙学观测预言}

\subsection{CMB 角功率谱修正}

TCSR 网络主要通过改变背景膨胀历史和晚期 ISW 效应影响 CMB。在大尺度 ($l < 30$) 上，温度功率谱 $C_l^{TT}$ 预计增强约 1–2\%，这是未来 CMB-S4 等实验可探测的特征信号。

\subsection{物质功率谱与大尺度结构}

TCSR 网络会导致物质功率谱在 $k \sim 0.1\, h/{\rm Mpc}$ 尺度处产生约 0.5–1\% 的压低，并轻微抑制结构增长因子。该效应在 Euclid 巡天的精度范围内有望被探测。

\subsection{哈勃常数 $H_0$ 与距离测量}

元宪模型预测 $H_0 \approx 67.4$ km/s/Mpc，与 Planck 结果一致，但与局部测量存在张力。该模型在 $z \sim 1.5$ 处与 ΛCDM 的光度距离差异约为 0.3\%。

\subsection{未来实验灵敏度预测}

Euclid + CMB-S4 联合分析有望在 $5$–$10$ 年内对 $w_a = 0.0012$ 的信号达到 $\geq 3\sigma$ 的检验能力。

\section{形式化验证与代码可用性}

\subsection{Lean 4 形式化验证}

所有核心推导已在 Lean 4 中完成机器验证，共包含 47 个主要定理，总代码行数约 2834 行，详见仓库 \texttt{lean/} 目录。

\subsection{可复现性保障}

本文所有数值结果均采用固定随机种子（seed=20260503）和高精度容差（rtol=$10^{-9}$），确保第三方可完全复现。

\subsection{数据与代码可用性}

完整代码仓库：  
\textbf{\url{https://github.com/YuanXian-Theory/YuanXian-Cosmology}}

所有资源采用 CC BY 4.0 协议开放获取。

\subsection{数据与代码可用性}

本仓库地址：  
\textbf{\url{https://github.com/YuanXian-Theory/YuanXian-Cosmology}}

\section{结论与展望}

本文在元宪理论框架下，完成了从微观 TCSR 量子场论到宏观修正弗里德曼方程的全链条形式化推导与数值实现，为探索暗能量的粒子物理本质提供了完备工具。

暗能量的本质，或许就藏在宇宙自指的深邃关联之中。

\begin{thebibliography}{20}

\bibitem{acharya2026tcsr}
Acharya, Z. (2026). Construction of the Complete Quantum Field Theory of the True Circle Self-Referon: Formal Derivation, Quantization, and Experimental Predictions in the YuanXian Framework. Zenodo. \doi{10.5281/zenodo.19942986}

\bibitem{acharya2026standard}
Acharya, Z. (2026). From Self-Consistent Universe to Computable Phenomena: Complete Formalization of the YuanXian Standard Model in Lean 4. Zenodo. \doi{10.5281/zenodo.19719792}

\bibitem{acharya2026yxtt}
Acharya, Z. (2026). YuanXian Self-Referential Mind Field Type Theory (YXTT 2.0): Mathematical Foundations, Physical Realization, and Unification of YuanXian Theory. Zenodo. \doi{10.5281/zenodo.19996191}

\bibitem{acharya2026axioms}
Acharya, Z. (2026). Formalization of the Four Core Axioms of YuanXian Theory and the ``All is One'' Unified Framework. Zenodo. \doi{10.5281/zenodo.19657426}

\bibitem{planck} Planck Collaboration. Planck 2018 results. VI. Cosmological parameters. A\&A, 2020.
\bibitem{desi} DESI Collaboration. The DESI Experiment Part I. arXiv:1611.00036, 2016.
\bibitem{cpl} Chevallier, M., \& Polarski, D. Accelerating universes with scaling dark matter. IJMPD, 2001.

\end{thebibliography}

\end{document}
